\documentclass{article}
%%%%%%%%%%%%%%%%%%%%%%%%%%%%%%%%%%%%%%%%%%%%%%%%%%%%%%%%%%%%%%%%%%%%%%%%%%%%%%%%%%%%%%%%%%%%%%%%%%%%%%%%%%%%%%%%%%%%%%%%%%%%%%%%%%%%%%%%%%%%%%%%%%%%%%%%%%%%%%%%%%%%%%%%%%%%%%%%%%%%%%%%%%%%%%%%%%%%%%%%%%%%%%%%%%%%%%%%%%%%%%%%%%%%%%%%%%%%%%%%%%%%%%%%%%%%
\usepackage{amsmath}
\usepackage{amssymb}
\usepackage{amsfonts}
\usepackage{xcolor}
\usepackage{booktabs}
\usepackage{caption}
\usepackage{subcaption}
%\usepackage{hyperref}
\usepackage{tikz}
\usepackage{ifthen} % needed for string comparisons below
\usepackage{xstring} 
\usepackage{listings}
\usepackage[capitalise]{cleveref}
\usepackage{bifurcationatlas_doc}
\bibliographystyle{chicagoa}
%\usepackage{multirow}
%\usepackage[usestackEOL]{stackengine}
\usepackage[paper=a4paper,margin=1in,left=2cm,right=2cm]{geometry}
% ----- marker shape mapping -----
%\def\shapeForblue{circle}
\graphicspath{{./fig/}}


\linespread{1.5}

\title{Bifurcation Encoding}
\author{Dieter Grass\thanks{TU Wien, Vienna, Austria; and International Institute for Applied Systems Analysis (IIASA), Laxenburg, Austria}}

\begin{document}
	
	\maketitle

\section{Latex package: Bifurcationatlas}
\subsection{Difference between \texttt{\textbackslash def} and \texttt{\textbackslash let}}

In TeX programming, \texttt{\textbackslash def} and \texttt{\textbackslash let} are both used to define macros, but they operate at fundamentally different levels.

\paragraph{\texttt{\textbackslash def}: macro definition via expansion}
The command \texttt{\textbackslash def} creates a new macro by specifying how it should expand when used. It performs textual replacement with full control over arguments.

\begin{lstlisting}[language=TeX]
\def\A{Hello}
\end{lstlisting}

Here, \texttt{\textbackslash A} expands to \texttt{Hello}. The definition is independent of any other macro and can include arguments:

\begin{lstlisting}[language=TeX]
\def\f#1{#1+1}
\end{lstlisting}

In this case, \texttt{\textbackslash f\{3\}} expands to \texttt{3+1}.

\paragraph{\texttt{\textbackslash let}: token-level assignment (aliasing)}
The command \texttt{\textbackslash let} assigns the current meaning of one control sequence to another. It does not redefine expansion logic but copies the existing internal meaning.

\begin{lstlisting}[language=TeX]
\let\B\A
\end{lstlisting}

This makes \texttt{\textbackslash B} behave exactly like \texttt{\textbackslash A} at the moment of assignment.

\paragraph{Key difference illustrated}
\begin{lstlisting}[language=TeX]
\def\A{Hello}
\let\B\A
\def\A{World}
\end{lstlisting}

After this sequence:
\begin{itemize}
  \item \texttt{\textbackslash A} expands to \texttt{World}
  \item \texttt{\textbackslash B} expands to \texttt{Hello}
\end{itemize}

This shows that \texttt{\textbackslash let} captures a \emph{snapshot} of a macro’s meaning, whereas \texttt{\textbackslash def} defines a \emph{re-evaluated expansion rule}.

\paragraph{Practical implication}
In macro-heavy systems (such as TikZ or key-value frameworks), \texttt{\textbackslash let} is useful for freezing state, while \texttt{\textbackslash def} is preferred for dynamic or data-dependent expansion.

\subsection{Key-Value Storage for Region Color Data}

The bifurcation atlas package uses the PGF key management system to decouple data specification from rendering. In particular, region-specific color information is stored via a dedicated key path:

\begin{lstlisting}[language=LaTeX]
\pgfkeys{
  /ba/region colors/.store in=\ba@colors
}
\end{lstlisting}

This declaration assigns a storage behavior to the key \texttt{/ba/region colors}. Whenever this key is set, its value is written into the macro \lstinline|\ba@colors|.

\paragraph{Mechanism.}
The directive \lstinline|.store in=\ba@colors| instructs PGF to treat \lstinline|\ba@colors| as a storage macro. When a user writes, for example:

\begin{lstlisting}[language=LaTeX]
\pgfkeys{/ba/region colors={red,blue,green}}
\end{lstlisting}

the internal PGF handler performs an assignment equivalent to:

\begin{lstlisting}[language=LaTeX]
\def\ba@colors{red,blue,green}
\end{lstlisting}

Thus, \lstinline|\ba@colors| acts as a container for the comma-separated list of colors associated with a given region.

\paragraph{Design rationale.}
Although \TeX{} macros are not variables in the classical programming sense, the use of \lstinline|\ba@colors| as a storage macro provides variable-like behavior within the expansion-based \TeX{} engine. The \texttt{@} character is used as a namespace separator to indicate that this macro is internal to the package and not part of the public interface. In standard \LaTeX{} conventions, such internal macros are protected from accidental user access unless \lstinline|\makeatletter| is enabled.

\paragraph{Summary.}
This mechanism enables a clean separation between user-level configuration (PGF keys) and internal data representation. The macro \lstinline|\ba@colors| serves as the central storage location for region color lists used in subsequent parsing and rendering stages of the bifurcation atlas system.

\subsection{Region Definition Interface}

Regions in the bifurcation atlas are defined through a dedicated user-level interface that binds a symbolic region name to a structured set of rendering data. This is implemented via the command:

\begin{lstlisting}[language=LaTeX]
\newcommand{\baNewRegion}[2]{%
% #1 = region name
% #2 = comma-separated color list

\tikzset{
  ba/region/#1/.style={
    /ba/region colors={#2}
  }
}%
}
\end{lstlisting}

\paragraph{Purpose.}
The macro \lstinline|\baNewRegion| provides a declarative mechanism for defining a region in the atlas. It associates a region identifier with a list of colors that will later be interpreted by the rendering system.

\paragraph{Arguments.}
\begin{itemize}
  \item \textbf{\#1 (region name):} A symbolic identifier used to reference the region in rendering commands.
  \item \textbf{\#2 (color list):} A comma-separated list of colors describing the region's internal structure.
\end{itemize}

\paragraph{Internal mechanism.}
The command expands into a TikZ style definition of the form:

\begin{lstlisting}[language=LaTeX]
ba/region/<name>/.style={
  /ba/region colors={...}
}
\end{lstlisting}

When this style is later applied, it sets the PGF key \texttt{/ba/region colors}, which stores the provided color list into the internal macro \lstinline|\ba@colors| via the key definition:

\begin{lstlisting}[language=LaTeX]
\pgfkeys{
  /ba/region colors/.store in=\ba@colors
}
\end{lstlisting}

\paragraph{Design rationale.}
This separation between region definition and storage allows a clean layering of abstraction: users interact only with named regions and color lists, while internal macros such as \lstinline|\ba@colors| remain hidden implementation details. The use of TikZ styles enables composability with other graphical transformations and ensures compatibility with the broader PGF/TikZ key-value system.

\subsection{Safe Access to Dynamically Stored Colors}

To safely retrieve colors that are generated dynamically during the parsing of a region definition, the package defines a guarded lookup macro:

\begin{lstlisting}[language=LaTeX]
\newcommand{\ba@getcolor}[1]{%
  \ifcsname ba@c#1\endcsname
    \csname ba@c#1\endcsname%
  \else
    white%
  \fi
}
\end{lstlisting}

\paragraph{Purpose.}
The macro \lstinline|\ba@getcolor| provides a robust interface for accessing internally stored color entries of the form \lstinline|\ba@c1|, \lstinline|\ba@c2|, etc. These entries are generated dynamically when a region's color list is processed.

\paragraph{Argument.}
\begin{itemize}
  \item \textbf{\#1:} Index of the color entry to retrieve (typically an integer such as 1--4 corresponding to the tiles of the bifurcation structure).
\end{itemize}

\paragraph{Mechanism.}
The macro uses the conditional construct

\begin{lstlisting}[language=LaTeX]
\ifcsname ba@c#1\endcsname ... \else ... \fi
\end{lstlisting}

to test whether the control sequence \lstinline|\ba@c#1| has been defined. If it exists, its value is expanded via:

\begin{lstlisting}[language=LaTeX]
\csname ba@c#1\endcsname
\end{lstlisting}

Otherwise, a safe fallback value \texttt{white} is returned.

\paragraph{Design rationale.}
Since the internal color storage is populated dynamically from comma-separated lists, the number of available entries may vary between regions. The guarded lookup prevents undefined control sequence errors and ensures that missing entries do not break rendering but instead default to a neutral color.

This design makes the rendering layer robust against incomplete or partially specified region data while preserving full flexibility in the underlying representation.

\subsection{Region Rendering Wrapper}

The rendering of a named region is encapsulated in a scalable TikZ wrapper that combines region lookup, graphical instantiation, and size control. This is implemented via:

\begin{lstlisting}[language=LaTeX]
\newcommand{\baRenderRegion}[2]{%
% #1 = region name
% #2 = scale factor

\scalebox{#2}{%
\begin{tikzpicture}
    \pic[ba/region/#1]{ba/bifstruct};
\end{tikzpicture}}%
}
\end{lstlisting}

\paragraph{Purpose.}
The macro \lstinline|\baRenderRegion| provides a unified interface for rendering a complete bifurcation region. It bridges the abstract region definition (identified by name) with its concrete graphical representation.

\paragraph{Arguments.}
\begin{itemize}
  \item \textbf{\#1 (region name):} The symbolic identifier of the region. This is resolved through the TikZ style system \texttt{ba/region/<name>}, which supplies the associated configuration data.
  \item \textbf{\#2 (scale factor):} A numeric scaling factor applied uniformly to the resulting TikZ picture.
\end{itemize}

\paragraph{Mechanism.}
The command constructs a standalone TikZ picture environment and invokes:

\begin{lstlisting}[language=LaTeX]
\pic[ba/region/#1]{ba/bifstruct};
\end{lstlisting}

This applies the region-specific style, which in turn initializes internal data structures such as color lists. The entire picture is then wrapped in a \lstinline|\scalebox| to allow flexible resizing without altering the underlying coordinate system.

\paragraph{Design rationale.}
This wrapper enforces a strict separation between geometry, data, and presentation. The region logic is handled entirely through TikZ styles, while scaling is handled externally via \lstinline|\scalebox|. This ensures that the same region definition can be reused consistently across different layouts (e.g., radial atlases or linear compositions) without modification.

\subsection{Radial Node Placement}

The radial layout system positions rendered regions on a circular structure using polar coordinates. This is implemented by the macro:

\begin{lstlisting}[language=LaTeX]
\newcommand{\baPlaceRadialNode}[4]{%
% #1 = index
% #2 = total number of nodes
% #3 = radius
% #4 = content

\pgfmathsetmacro{\ba@angle}{90-(#1-1)*360/#2}

\node at (\ba@angle:#3) {#4};
}
\end{lstlisting}

\paragraph{Purpose.}
The macro \lstinline|\baPlaceRadialNode| computes angular positions for nodes in a circular arrangement and places arbitrary content at the corresponding polar coordinates.

\paragraph{Arguments.}
\begin{itemize}
  \item \textbf{\#1 (index):} The position of the node in the sequence (starting from 1).
  \item \textbf{\#2 (total number of nodes):} The total number of elements in the radial layout, used to compute equal angular spacing.
  \item \textbf{\#3 (radius):} The distance from the origin at which the node is placed.
  \item \textbf{\#4 (content):} The graphical content to be placed at the computed position, typically a rendered region.
\end{itemize}

\paragraph{Angle computation.}
The angular position is computed using:

\begin{lstlisting}[language=LaTeX]
\pgfmathsetmacro{\ba@angle}{90-(#1-1)*360/#2}
\end{lstlisting}

This formula distributes nodes evenly along a circle, starting from the top (90 degrees) and proceeding clockwise. The subtraction of 90 degrees ensures a conventional orientation where the first element appears at the top of the circle.

\paragraph{Placement mechanism.}
The computed angle is then used in TikZ polar coordinate syntax:

\begin{lstlisting}[language=LaTeX]
\node at (\ba@angle:#3) {#4};
\end{lstlisting}

which places the node at radius \lstinline|#3| and angle \lstinline|\ba@angle|.

\paragraph{Design rationale.}
This abstraction decouples geometric computation from rendering logic. It allows higher-level layout functions (such as radial atlases) to operate purely in terms of indexed regions without managing explicit coordinate calculations. The use of PGF mathematical evaluation ensures consistency and numerical accuracy across all layout operations.

\subsection{Sequential Composition of Regions}

The simplest assembly mechanism in the bifurcation atlas is a linear composition of multiple regions. This is implemented by iterating over a comma-separated list of region identifiers and rendering each one sequentially.

\begin{lstlisting}[language=LaTeX]
\newcommand{\baComposeRegions}[2]{%
% #1 = comma-separated region list
% #2 = scale factor

\foreach \r in {#1}{%
    \baRenderRegion{\r}{#2}%
}%
}
\end{lstlisting}

\paragraph{Purpose.}
The macro \lstinline|\baComposeRegions| provides a lightweight mechanism for rendering multiple regions in sequence. It does not impose any geometric structure; instead, it delegates layout entirely to the surrounding TikZ context.

\paragraph{Arguments.}
\begin{itemize}
  \item \textbf{\#1 (region list):} A comma-separated list of region identifiers to be rendered in order.
  \item \textbf{\#2 (scale factor):} A uniform scaling factor applied to each region via \lstinline|\baRenderRegion|.
\end{itemize}

\paragraph{Mechanism.}
The implementation uses the PGF iteration construct:

\begin{lstlisting}[language=LaTeX]
\foreach \r in {#1}{ ... }
\end{lstlisting}

which expands the region list and assigns each entry in turn to the macro \lstinline|\r|. Each region is then passed to the rendering layer:

\begin{lstlisting}[language=LaTeX]
\baRenderRegion{\r}{#2}
\end{lstlisting}

ensuring consistent visualization of all elements.

\paragraph{Design rationale.}
This composition layer is intentionally minimalistic. It separates iteration logic from rendering logic, allowing the same rendering function to be reused in more complex layouts (such as radial or hierarchical arrangements). It serves as a baseline assembly mechanism for simple linear visualization of multiple bifurcation regions.
\end{document}